\documentclass[12pt,a4paper]{article}

% Encoding and language
\usepackage[T1]{fontenc}
\usepackage[utf8]{inputenc}
\usepackage[english]{babel}

% Page layout
\usepackage[margin=2.5cm]{geometry}

% Mathematics
\usepackage{amsmath,amssymb}

% Better code formatting
\usepackage{listings}

\title{Java and Spring Boot Interview Notes}
\author{}
\date{\today}

\begin{document}

\maketitle

\section{Java Interview Questions}

\subsection{What is the Difference Between \texttt{==} and \texttt{equals()}?}

In Java, \texttt{==} compares primitive values directly. For objects, it
checks whether two references point to the same object. The
\texttt{equals()} method is used for logical equality and can be overridden
by a class to compare object contents. The default implementation of
\texttt{equals()} in the \texttt{Object} class performs reference
comparison.

\subsubsection*{Key Point}

\begin{itemize}
    \item For primitives, \texttt{==} compares values.
    \item For objects, \texttt{==} compares references.
    \item \texttt{equals()} is used for logical or content equality.
\end{itemize}

\subsection{What Does the \texttt{final} Keyword Mean for Variables, Methods, and Classes?}

In Java, the \texttt{final} keyword prevents further modification depending
on where it is used. A final variable cannot be reassigned, a final method
cannot be overridden and a final class cannot be extended. For object
references, \texttt{final} prevents reassignment of the reference, but it
does not make the object itself immutable.

\subsection{What is the \texttt{@Qualifier} Annotation in Spring?}

The \texttt{@Qualifier} annotation is used when multiple beans of the same
type exist in the Spring ApplicationContext and Spring cannot determine
which one should be injected.

\subsubsection*{\texttt{@Qualifier} vs \texttt{@Primary}}

\texttt{@Primary} defines the default bean when multiple beans of the same
type exist.

\texttt{@Qualifier} explicitly selects a specific bean.

\begin{verbatim}
@Primary   -> preferred/default bean
@Qualifier -> explicitly selected bean
\end{verbatim}

\subsubsection*{Important Note}

Using \texttt{@Qualifier} binds that dependency to a specific implementation.

Therefore, it is suitable when the implementation is known when the
application is configured.

If the implementation must be selected dynamically at runtime, for example
based on a user's payment method, a Strategy or Factory-based design is
usually more appropriate.

\subsubsection*{Interview Answer}

The \texttt{@Qualifier} annotation is used when multiple beans of the same
type exist in the Spring context. It tells Spring exactly which bean should
be injected and resolves ambiguity during dependency injection.

It is mainly used for static dependency selection. If the implementation
must be chosen dynamically at runtime, a Strategy or Factory approach can
be used instead.

\subsection{Can a Static Method Be Overridden in Java?}

No. Static methods cannot be overridden in Java.

If a subclass defines a static method with the same signature as a static
method in its superclass, this is called \textit{method hiding}, not method
overriding.

\subsubsection*{Method Overriding vs Method Hiding}

\begin{verbatim}
Instance method:
-> Can be overridden
-> Runtime polymorphism
-> Actual object type determines the method

Static method:
-> Cannot be overridden
-> Can be hidden
-> Reference/class type determines the method
\end{verbatim}

\subsubsection*{Interview Answer}

Static methods cannot be overridden in Java because they belong to the
class rather than an instance. If a subclass defines a static method with
the same signature, it hides the superclass method. This is called method
hiding. Static method calls are resolved based on the reference or class
type rather than the actual object type.

\subsection{Which Memory Area is Cleaned by the Garbage Collector: Heap or Stack?}

The Java Garbage Collector primarily manages \textbf{heap memory}.

Objects created with \texttt{new} are generally allocated on the heap.

For example:

\begin{lstlisting}[language=Java]
Person person = new Person("Utku");
\end{lstlisting}

The local reference \texttt{person} may exist in the current stack frame,
while the actual \texttt{Person} object is stored on the heap.

If the object becomes unreachable, it becomes eligible for garbage
collection.

\begin{lstlisting}[language=Java]
Person person = new Person("Utku");

person = null;
\end{lstlisting}

After the reference is removed, if no other reachable reference points to
the object, the object may be reclaimed by the Garbage Collector.

Being eligible for garbage collection does not mean that the object is
immediately removed. The JVM decides when garbage collection should run.

\subsubsection*{What Happens to Stack Memory?}

Stack memory is managed automatically using stack frames.

Each method call creates a stack frame that contains information such as
local variables and references.

\begin{lstlisting}[language=Java]
void calculate() {
    int x = 10;
    int y = 20;
}
\end{lstlisting}

When the method finishes, its stack frame is automatically removed.
Therefore, the Garbage Collector does not normally clean stack memory.

\begin{verbatim}
Heap:
-> Stores objects
-> Managed mainly by the Garbage Collector
-> Unreachable objects may be reclaimed

Stack:
-> Stores method frames and local variables
-> Automatically managed
-> Frames are removed when methods return
\end{verbatim}

\subsubsection*{Interview Answer}

The Java Garbage Collector mainly manages heap memory. Objects are allocated
on the heap and when they are no longer reachable, they become eligible for
garbage collection. Stack memory is managed automatically through stack
frames, which are removed when methods return.

\subsection{Why is \texttt{String} Immutable in Java?}

A \texttt{String} object cannot be modified after it has been created.
Operations that appear to modify a String actually create a new String
object.

\subsubsection*{Interview Answer}

String is immutable in Java mainly for security, safe String Pool sharing,
thread safety, and reliable hashing. Since a String cannot change after
creation, it can safely be shared between threads, reused from the String
Pool, and used as a stable key in hash-based collections such as HashMap.

\subsection{What is the Difference Between \texttt{HashSet} and \texttt{TreeSet}?}

Both \texttt{HashSet} and \texttt{TreeSet} implement the \texttt{Set}
interface, meaning that they do not allow duplicate elements.

\begin{verbatim}
HashSet:
-> Unique elements
-> No guaranteed ordering
-> Hash-based
-> Uses hashCode() and equals()
-> Average O(1) operations

TreeSet:
-> Unique elements
-> Sorted elements
-> Tree-based
-> Uses Comparable or Comparator
-> O(log n) operations
\end{verbatim}

\subsubsection*{Interview Answer}

Both HashSet and TreeSet store unique elements. HashSet uses hashing and
does not guarantee ordering, while TreeSet keeps elements sorted using a
tree structure. HashSet provides average O(1) operations, whereas TreeSet
provides O(log n) operations. HashSet relies mainly on hashCode() and
equals(), while TreeSet relies on Comparable or Comparator.

\subsection{What is the Difference Between \texttt{HashMap} and \texttt{HashSet}?}

\texttt{HashMap} stores data as key-value pairs, while \texttt{HashSet}
stores only unique elements.

\subsubsection*{Interview Answer}

HashMap stores key-value pairs, while HashSet stores only unique elements.
HashMap does not allow duplicate keys, although duplicate values are
allowed. HashSet does not allow duplicate elements. Internally, HashSet is
backed by a HashMap, where the set elements are stored as map keys with a
dummy value.

\subsection{What is a Bean?}

A bean is simply an object whose lifecycle is managed by Spring.

Spring may create beans automatically by discovering classes annotated
with annotations such as:

\begin{verbatim}
@Component
@Service
@Repository
@Controller
@RestController
\end{verbatim}

Beans can also be explicitly declared using the \texttt{@Bean}
annotation.

\subsection{What is the ApplicationContext?}

The \texttt{ApplicationContext} is the central Spring IoC container.

It is responsible for creating, storing, configuring, and connecting
Spring-managed objects, which are called \textbf{beans}.

\texttt{ApplicationContext} can be thought of as the
container in which Spring-managed objects live.

\subsection{What Does \texttt{@SpringBootApplication} Do Internally?}

\texttt{@SpringBootApplication} is a convenience annotation that combines:

\begin{verbatim}
@SpringBootConfiguration
@EnableAutoConfiguration
@ComponentScan
\end{verbatim}

\texttt{@SpringBootConfiguration} is based on
\texttt{@Configuration} and allows Spring configuration and bean
definitions.

For example:

\begin{lstlisting}[language=Java]
@Configuration
public class AppConfig {

    @Bean
    public PaymentClient paymentClient() {
        return new PaymentClient();
    }
}
\end{lstlisting}

The object returned by a \texttt{@Bean} method is registered in the
Spring \texttt{ApplicationContext}.

\texttt{@ComponentScan} searches the application's package and its
subpackages for classes such as:

\begin{verbatim}
@Component
@Service
@Repository
@Controller
@RestController
\end{verbatim}

and registers them as Spring beans.

\texttt{@EnableAutoConfiguration} automatically configures the application
based on the libraries available on the classpath, application properties
and existing beans.

For example, if \texttt{spring-boot-starter-web} is present, Spring Boot
automatically configures much of the Spring MVC infrastructure.

Finally:

\begin{lstlisting}[language=Java]
SpringApplication.run(DemoApplication.class, args);
\end{lstlisting}

starts the application and initializes the Spring
\texttt{ApplicationContext}.

\subsubsection*{Interview Answer}

\texttt{@SpringBootApplication} combines configuration, component scanning
and auto-configuration. It allows Spring to discover and create beans,
configure the application automatically and manage those beans inside the
\texttt{ApplicationContext}. The \texttt{ApplicationContext} is the IoC
container responsible for managing beans and injecting their dependencies.

\subsection{What is \texttt{@Autowired} in Spring?}

\texttt{@Autowired} is used for dependency injection. It tells Spring to
find a matching bean from the \texttt{ApplicationContext} and inject it.

\begin{lstlisting}[language=Java]
@Service
public class OrderService {

    private final PaymentService paymentService;

    public OrderService(PaymentService paymentService) {
        this.paymentService = paymentService;
    }
}
\end{lstlisting}

Since Spring 4.3, if a class has only one constructor,
\texttt{@Autowired} is not required because Spring automatically uses that
constructor for dependency injection.

If a class has multiple constructors, \texttt{@Autowired} can be used to
indicate which constructor Spring should use.

\texttt{final} is commonly used with constructor injection because the
dependency is assigned once during object creation and cannot later be
reassigned to another object.

\texttt{@Autowired} can also be used for setter and field injection, but
constructor injection is generally preferred.

\subsubsection*{Interview Answer}

\texttt{@Autowired} tells Spring to inject a matching bean from the
\texttt{ApplicationContext}. With a single constructor, it is usually not
needed. With multiple constructors, it can specify which constructor should
be used for injection. Constructor-injected dependencies are often declared
\texttt{final} so their references cannot be reassigned.
\end{document}